\documentclass[11pt]{article}

\usepackage[a4paper,margin=24mm]{geometry}
\usepackage[T1]{fontenc}
\usepackage[utf8]{inputenc}
\usepackage{amsmath,amssymb}
\usepackage{array}
\usepackage{booktabs}
\usepackage{enumitem}
\usepackage{longtable}
\usepackage{tabularx}
\usepackage{xcolor}
\usepackage{listings}
\usepackage[hidelinks]{hyperref}

\setlength{\parindent}{0pt}
\setlength{\parskip}{0.55em}
\setlist{leftmargin=*,nosep}
\emergencystretch=3em
\renewcommand{\arraystretch}{1.16}

\definecolor{codebg}{RGB}{247,247,247}
\definecolor{rulegrey}{RGB}{92,92,92}
\lstset{
  basicstyle=\ttfamily\small,
  backgroundcolor=\color{codebg},
  frame=single,
  rulecolor=\color{rulegrey},
  breaklines=true,
  columns=fullflexible,
  keepspaces=true,
  showstringspaces=false,
  aboveskip=0.7em,
  belowskip=0.7em
}

\DeclareRobustCommand{\code}[1]{\texttt{\detokenize{#1}}}
\newcommand{\ABALearn}{\textsc{ABALearn}}
\newcommand{\Hprobe}[1]{\textbf{H#1}}
\newcommand{\Epos}{E^{+}}
\newcommand{\Eneg}{E^{-}}
\newcommand{\notindep}{\not\!\perp\!\!\!\perp}
\newcommand{\indep}{\perp\!\!\!\perp}
\newcommand{\evidence}{\paragraph{Evidence.}}
\newcommand{\boundary}{\paragraph{Boundary.}}
\newcommand{\interpretation}{\paragraph{Interpretation.}}

\title{What Target-wise ABA Learning Learns from\\
Deterministic Causal Fixtures\\[0.35em]
\large M1.3 Bucket 3 synthesis of H0--H7b}
\author{Working research synthesis for discussion with Fabrizio}
\date{4 August 2026}

\begin{document}
\maketitle

\begin{center}
\fbox{\begin{minipage}{0.92\textwidth}
\textbf{Document status.}
This is a supervisor-facing working synthesis of the completed H0--H7b
evidence. It is not submitted-report prose and it does not lock a Bucket~3
claim. H7a and H7b are bounded implementation-behaviour probes of assumption
introduction under brave and cautious nd respectively; neither is a method
ranking or a causal-recovery result.
The document distinguishes direct observations, trace-supported explanations,
bounded interpretations, and candidate claims for discussion.
\end{minipage}}
\end{center}

\begin{abstract}
This investigation asks what target-wise ABA Learning does when its only
learner-visible input is a finite binary table sampled from a known causal
fixture. The central fixture is the deterministic AND collider
\(A\rightarrow C\leftarrow B\), with independent stochastic roots and
\(C=A\land B\). The graph, population, finite table, learner encoding, learned
framework, and evaluator-only mechanism reference are kept separate. The
desirable deterministic target is \(C\); root-target runs are retained as
diagnostics of predictive overreach and search behaviour.

Across H0--H7b, \ABALearn{} sometimes emits the direct AND rule, but this
result is not invariant to sample support, irrelevant variables, background
predicate order, search strategy, learning semantics, folding budget, or
assumption-introduction policy. A procedural \code{solved} outcome can denote
a mechanism-aligned rule, a non-minimal distractor-retaining theory, or a brave
assumption construction that does not define the target from its observed
predictors. Under brave nd, \code{sechk} produces a first-BK latch that drops
later mechanism parents on the tested child targets; under cautious nd it can
reach those parents later through deeper assumption nests. The experiments
therefore provide a controlled account of
\emph{how} and \emph{why} target-wise ABA Learning succeeds or departs from
the evaluator mechanism on these fixtures; they do not yet constitute
Russo-style Causal ABA, CPDAG recovery, or a graph decoder.
\end{abstract}

\tableofcontents
\newpage

\section*{Executive findings}
\addcontentsline{toc}{section}{Executive findings}

The completed evidence supports the following bounded synthesis.

\begin{enumerate}
  \item \textbf{Exact mechanism-rule recovery is possible but not invariant.}
  On the baseline sample, Greedy AAMAS learning of target \(c\) emits the
  direct conjunction:
\begin{lstlisting}
c(A) :- a_val_1(A), b_val_1(A).
\end{lstlisting}
  Nd learning solves the same target with assumption-rich representations,
  while added irrelevant variables are retained by both Greedy and nd.

  \item \textbf{\code{solved} is not a scientific correctness criterion.}
  It records successful completion under the selected semantics. It does not
  establish minimality, mechanism agreement, or causal graph recovery.
  Counterexamples include brave nd root solutions and solved child theories
  containing isolated variables.

  \item \textbf{Finite support can reverse apparent root learnability.}
  Removing the rare \((0,0,0)\) witness makes Greedy emit false both-zero
  rules for roots \(a\) and \(b\); restoring it makes the same rules fail.

  \item \textbf{Irrelevant variables and background order affect structure.}
  Greedy maximal co-occurrence retains an isolated variable in rule bodies.
  Nd can select a BK-leading isolated variable at the first fold and build the
  final framework around it.

  \item \textbf{Brave versus cautious effects depend on the search bundle.}
  Under nd/\code{relto}, cautious checking blocks a brave residual-assumption
  gadget and changes outcomes or a contrary closure. Under
  Greedy/\code{relto}, the semantic flip is inert on all 18 tested pairs,
  apart from runtime.

  \item \textbf{Long cautious-nd failures are repeated search.}
  Increasing \code{folding_steps} replays essentially the same failed root
  motif; increasing \(n\) enlarges per-band grounding and bookkeeping.

  \item \textbf{Assumption introduction interacts with learning semantics.}
  Under brave nd, changing \code{relto} to \code{sechk} changes all seven
  paired deltas: \code{sechk} latches onto the first BK predictor and omits
  later covariates. Under cautious nd, the two comparable child deltas also
  differ, but \code{sechk} reaches the mechanism parents through deeper nests
  and incurs more search. The brave pattern therefore does not transfer
  unchanged to cautious learning.
\end{enumerate}

These findings describe the tested fixtures and configurations. They are not
universal properties of ABA Learning and do not yet establish a Bucket~3
claim.

\section{Aim, scope, and interpretative standard}

\subsection{Research question}

\begin{quote}
Given a finite table sampled from a known deterministic causal fixture, what
rules and ABA structures does target-wise \ABALearn{} produce, under which
conditions do those structures correspond to the generating mechanism, and
what explains departures from it?
\end{quote}

The investigation begins below full causal graph discovery. A frozen table is
encoded separately for each target. Target value \(1\) supplies \(\Epos\);
target value \(0\) supplies \(\Eneg\); exact values of all non-target variables
form the BK predicates. The learner is not told the graph, causal roles,
population, mechanism reference, or CPDAG.

All variables are run as targets to preserve target agnosticism at execution
level. Evaluation remains role-sensitive: deterministic child \(c\) has a
mechanism to recover, while a stochastic root has no deterministic rule over
the other observed variables.

\subsection{Objects kept separate}

\begin{enumerate}
  \item generating DAG \(G\);
  \item structural mechanisms;
  \item exact observational population \(P\);
  \item finite IID sample \(\mathcal{D}_n\sim P\);
  \item learner-visible BK and examples derived from \(\mathcal{D}_n\);
  \item learned ABA framework and delta;
  \item procedural outcome under brave or cautious semantics;
  \item evaluator-only mechanism reference; and
  \item graph-level recovery object such as a DAG, MEC, or CPDAG.
\end{enumerate}

A rule may cover the sample without being a population implication; a
framework may satisfy brave examples without defining a function from BK to
target; and a rule may mention a parent without orienting a causal edge.

\subsection{Evaluation vocabulary}

\begin{center}
\begin{tabularx}{\textwidth}{>{\raggedright\arraybackslash}p{0.24\textwidth}X}
\toprule
Term & Meaning in this document \\
\midrule
\code{solved} & Accepted serialized framework under the selected semantics; a
procedural outcome. \\
No solution & The recorded \code{completed_no_solution} outcome: search
terminated without an accepted theory, relative to the configured strategy
and budget. \\
\code{timeout} & Wall-time limit reached; not evidence of unlearnability. \\
Sample-adequate & Framework meets encoded
\(\langle\Epos,\Eneg\rangle\) on observed rows under the selected semantics. \\
Mechanism-aligned string & Learned rule is syntactically identical to the
evaluator rule; stronger than \code{solved}, but not graph recovery. \\
Distractor-retaining & Learned theory refers to a variable absent from the
target mechanism. \\
Spurious finite-sample rule & Learned implication holds in the sample but is
false in the population. \\
Causal or graph recovery & Conclusion about mechanism, edges, MEC, or CPDAG;
none follows automatically from the preceding categories. \\
\bottomrule
\end{tabularx}
\end{center}

Satisfiability of \code{bk.sol_chk.asp} is a final-artefact integrity witness,
not an additional recovery metric and, under cautious learning, not proof that
every positive is a cautious consequence.

\subsection{Relation to Causal ABA}

The current pipeline is target-wise ABA Learning over tabular data. It learns
rules, assumptions, and contraries by transforming BK and examples. It does
not implement the Causal ABA machinery in which assumptions encode arrows,
no-edge statements, or conditional independences, rules encode acyclicity and
\(d\)-separation, and stable extensions correspond to candidate DAGs. It
therefore should not be described as full causal discovery.

\section{Causal fixtures and evaluator ground truth}

\subsection{Baseline deterministic AND collider}

\[
G_0:\qquad A\longrightarrow C\longleftarrow B,
\]
with learner identifiers \(a,b,c\), and
\begin{align*}
A &\sim \operatorname{Bernoulli}(4/5),&
B &\sim \operatorname{Bernoulli}(7/10),\\
A &\indep B,& C&=A\land B.
\end{align*}

\begin{center}
\begin{tabular}{@{}cccc@{}}
\toprule
\(A\) & \(B\) & \(C\) & \(P(A,B,C)\) \\
\midrule
0 & 0 & 0 & \(3/50=0.06\) \\
0 & 1 & 0 & \(7/50=0.14\) \\
1 & 0 & 0 & \(6/25=0.24\) \\
1 & 1 & 1 & \(14/25=0.56\) \\
\bottomrule
\end{tabular}
\end{center}

The other four binary assignments are structural zeros. Causal sufficiency is
declared through mutually independent root exogenous noise. The Markov
condition is verified by
\[
P(a,b,c)=P(a)P(b)P(c\mid a,b).
\]
Ordinary faithfulness is verified exactly by an exhaustive singleton-pair
audit. The only population independence is \(A\indep B\); in particular,
\(A\notindep B\mid C\), and each parent is dependent on \(C\).

The ordinary CI MEC has one member and its CPDAG is
\(A\rightarrow C\leftarrow B\). This is evaluator-side population truth.
Deterministic relations can require recovery objects beyond ordinary CI, so no
stronger deterministic-identifiability result is asserted.

\subsection{Evaluator-only mechanism reference}

\begin{lstlisting}
c(A) :- a_val_1(A), b_val_1(A).
\end{lstlisting}

Targets \(a\) and \(b\) have
\code{no_observed_parent_deterministic_rule}: they are stochastic roots. The
graph, certificate, and mechanism reference are never learner-visible.

\subsection{Frozen baseline sample}

The seed-42 IID sample with \(n=30\) has counts:

\begin{center}
\begin{tabular}{@{}lr@{}}
\toprule
Assignment & Count \\
\midrule
\((1,1,1)\) & 17 \\
\((1,0,0)\) & 7 \\
\((0,1,0)\) & 4 \\
\((0,0,0)\) & 2 \\
\bottomrule
\end{tabular}
\end{center}

All population atoms appear. Example counts are \(a:24/6\), \(b:21/9\), and
\(c:17/13\), written as \(|\Epos|/|\Eneg|\). The same table is used for all
H0 targets and configurations.

For \(c\), all positives share \(A=B=1\). For target \(a\), predictor cell
\((B,C)=(0,0)\) contains both labels: row 9 is \((1,0,0)\), while row 26 is
\((0,0,0)\). No deterministic function of \((B,C)\) distinguishes them.
Target \(b\) has the symmetric ambiguity.

\subsection{Controlled fixture extensions}

\begin{center}
\begin{tabularx}{\textwidth}{
  >{\raggedright\arraybackslash}p{0.22\textwidth}
  >{\raggedright\arraybackslash}p{0.25\textwidth}X}
\toprule
Fixture & Structure & Purpose \\
\midrule
Baseline & \(A\rightarrow C\leftarrow B\), \(C=A\land B\) &
Target geometry and AAMAS/ECAI baseline. \\
Isolated trailing \(D\) & Same collider and AND; isolated
\(D\sim\operatorname{Bernoulli}(1/2)\) &
Test Greedy retention of an irrelevant trailing BK block. \\
Isolated leading \(A\) & \(B\rightarrow C\leftarrow D\), \(C=B\land D\);
isolated \(A\sim\operatorname{Bernoulli}(1/2)\) &
Place an irrelevant variable first in target-\(c\) BK order \((a,b,d)\). \\
\bottomrule
\end{tabularx}
\end{center}

Each extension has verified Markov factorisation and ordinary faithfulness.
Evaluator rules contain only the generating parents:

\begin{lstlisting}
% Collider plus isolated D:
c(A) :- a_val_1(A), b_val_1(A).
% BD collider plus isolated A:
c(A) :- b_val_1(A), d_val_1(A).
\end{lstlisting}

\section{Experimental logic and comparison strength}

\subsection{Learner configurations}

\begin{center}
\small
\begin{tabularx}{\textwidth}{
  >{\raggedright\arraybackslash}p{0.21\textwidth}
  >{\raggedright\arraybackslash}p{0.10\textwidth}
  >{\raggedright\arraybackslash}p{0.14\textwidth}
  >{\raggedright\arraybackslash}p{0.12\textwidth}
  >{\raggedright\arraybackslash}p{0.12\textwidth}X}
\toprule
Identity & Semantics & Folding & Selection & Space & Assumption introduction \\
\midrule
\code{aamas2025} & brave & greedy & mgr & BK & \code{relto} \\
\code{ecai2024} & brave & nd & any & all & \code{relto} \\
\code{baseline_cautious} & cautious & nd & any & all & \code{relto} \\
\code{greedy_cautious} & cautious & greedy & mgr & BK & \code{relto} \\
\code{nd_brave_sechk} & brave & nd & any & all & \code{sechk} \\
\code{nd_cautious_sechk} & cautious & nd & any & all & \code{sechk} \\
\bottomrule
\end{tabularx}
\end{center}

\code{greedy_cautious}, \code{nd_brave_sechk}, and
\code{nd_cautious_sechk} are experimental, not published AAMAS or ECAI
methods. H6 configurations match
\code{baseline_cautious} except for \code{folding_steps}.

\subsection{Probe map}

\begin{longtable}{@{}p{0.075\textwidth}p{0.24\textwidth}p{0.29\textwidth}p{0.31\textwidth}@{}}
\toprule
Probe & Question & Intervention & Principal evidence \\
\midrule
\endfirsthead
\toprule
Probe & Question & Intervention & Principal evidence \\
\midrule
\endhead
\Hprobe{0} & Baseline behaviour by target & Same fixture/table; two strategy
bundles & Six cells and traces. \\
\Hprobe{1} & Missing support witness & Nested \(n=25\) versus \(n=30\), same
seed/config & Root deltas and rows 26/30. \\
\Hprobe{2} & Trailing irrelevant covariate & Add isolated \(D\); AAMAS target
\(c\) & Two \(D\)-split conjunctions. \\
\Hprobe{3} & BK-leading distractor & BD AND plus isolated leading \(A\) &
First fold, \(a\)-gated delta, AAMAS contrast. \\
\Hprobe{4} & Brave versus cautious on roots & Semantics only under
nd/\code{relto} & Root outcomes and residual closure. \\
\Hprobe{4b} & Brave versus cautious on H3 \(c\) & Same H3 table and nd bundle &
One changed \code{c_alpha_3} rule. \\
\Hprobe{5} & Brave versus cautious in Greedy & Semantics only; 18 pairs &
Outcomes, byte-identical deltas, traces, runtime. \\
\Hprobe{6a} & Folding-budget cost & \(M\in\{1,2,5,10\}\), fixed \(n=30\) &
Token replays, lines, runtime. \\
\Hprobe{6b} & Nested-\(n\) cost & \(n\in\{30,60,90\}\), fixed \(M=2\) &
Fixed topology, growing per-band work. \\
\Hprobe{7a} & \code{relto} versus \code{sechk} & Assumption policy only under
brave nd & Seven pairs and trace forks. \\
\Hprobe{7b} & \code{relto} versus \code{sechk} under cautious nd & Assumption
policy only on H0/H3 & Child deltas, failed-root cost, timeout limits. \\
\bottomrule
\end{longtable}

\subsection{Three levels of comparison}

\paragraph{Level 1: exact paired option ablations.}
H4, H4b, H5, H6a, H7a, and H7b hold fixture, sample, encoding, and relevant inputs
fixed while changing one option. They support the strongest procedural
attribution.

\paragraph{Level 2: controlled data or fixture interventions.}
H1, H2, H3, and H6b alter the sample or fixture. H1 changes \(n\) and observed
support; H6b changes multiplicities and computational size; H2/H3 change the
observed variable set and order.

\paragraph{Level 3: bundle-level qualitative contrasts.}
AAMAS versus ECAI and Greedy-cautious versus baseline cautious change several
search settings. They show behavioural differences but cannot attribute them
to one option.

\section{Baseline behaviour: target geometry and search strategy}

\subsection{Outcome overview}

\begin{center}
\small
\begin{tabularx}{\textwidth}{
  >{\raggedright\arraybackslash}p{0.12\textwidth}
  >{\raggedright\arraybackslash}p{0.10\textwidth}
  >{\raggedright\arraybackslash}p{0.20\textwidth}X}
\toprule
Arm & Target & Outcome & Learned form or failure \\
\midrule
AAMAS & \(c\) & \code{solved} & One assumption-free AND conjunction,
string-identical to the evaluator rule. \\
AAMAS & \(a,b\) & no solution & Both-one rule safe;
both-zero rule conflicts with rows 26/30 and repair dead-ends. \\
ECAI & \(c\) & \code{solved} & Under-fold to \(a=1\), add \(\alpha_1\),
attack it when \(b=0\). \\
ECAI & \(a,b\) & \code{solved} & Brave per-row choice in an ambiguous
predictor cell; no root mechanism recovered. \\
\bottomrule
\end{tabularx}
\end{center}

\subsection{Greedy AAMAS on the child}

Every target-\(c\) positive lies in \(A=B=1\). Greedy BK folding constructs
maximal co-occurrence body
\(\{\code{a_val_1},\code{b_val_1}\}\); mgr removes duplicates. No negative
shares it, so no assumption is introduced:

\begin{lstlisting}
c(A) :- a_val_1(A), b_val_1(A).
\end{lstlisting}

This syntactically matches the evaluator rule. It is the strongest H0
mechanism-alignment observation, but remains bounded to target, table,
encoding, and configuration. It does not recover the graph or establish
irrelevant-variable robustness.

\subsection{Greedy AAMAS on the roots}

For target \(a\), positives occupy \((B,C)=(1,1)\) and \((0,0)\). Greedy
proposes:

\begin{lstlisting}
a(A) :- b_val_1(A), c_val_1(A).
a(A) :- b_val_0(A), c_val_0(A).
\end{lstlisting}

The first is safe. The second covers negative rows 26 and 30. Assumption
introduction creates a defeasible both-zero rule, but contrary folding
reconstructs the same body. Since the assumption is already present,
\code{relto} cannot introduce another; the trace ends KO and
\code{completed_no_solution}. Target \(b\) is symmetric.

This is evaluator-compatible only in that no false deterministic root rule is
emitted. It is not root-mechanism recovery.

\subsection{Nd ECAI on the child}

Nd folds one predicate at a time. With order \(a,b\), the first fold yields
\(\code{c(A) :- a_val_1(A).}\), which covers seven \(A=1,B=0\) negatives.
Assumption introduction follows; after an unsuccessful \(a\)-based contrary
fold, \code{relto} reaches \code{b_val_0}:

\begin{lstlisting}
c(A) :- alpha_1(A), a_val_1(A).
c_alpha_1(A) :- b_val_0(A).
assumption(alpha_1(A)).
contrary(alpha_1(A),c_alpha_1(A)) :- assumption(alpha_1(A)).
\end{lstlisting}

On this binary table, the theory derives \(c\) when \(a=1\) unless \(b=0\)
attacks the required assumption. It is brave sample-adequate and mentions both
parents across target and contrary rules, but is not the canonical conjunction
and is not recorded as mechanism recovery.

\subsection{Nd ECAI on the roots: brave choice construction}

For target \(a\), ECAI first constructs a safe \(b=1\) block guarded by an
assumption attacked at \(c=0\). The ambiguous \((B,C)=(0,0)\) cell then
produces:

\begin{lstlisting}
a(A) :- alpha_2(A), b_val_0(A).
c_alpha_2(A) :- alpha_3(A), c_val_0(A).
c_alpha_3(A) :- alpha_2(A), b_val_0(A).
\end{lstlisting}

Rows 9 and 26 have identical BK but opposite labels. The ungrounded rules do
not define \(a\) as a function of \((b,c)\). Different ground assumptions can
take different statuses in a witnessing stable model. Brave learning requires
existence of a stable model satisfying the joint task, not a unique
predictor-to-target mapping.

\interpretation
H0 exposes distinct notions of success. AAMAS \(c\) produces a
mechanism-aligned string. ECAI roots pass the brave task despite predictor
ambiguity. Treating both merely as \code{solved} erases that distinction.

\section{Finite-sample support: H1}

H1 uses the first 25 rows of the same seed-42 stream and AAMAS configuration.
The prefix contains \((1,1,1):14\), \((1,0,0):7\), and \((0,1,0):4\), with no
\((0,0,0)\). The population still gives that atom probability \(3/50\).

\begin{center}
\begin{tabularx}{\textwidth}{
  >{\raggedright\arraybackslash}p{0.11\textwidth}
  >{\raggedright\arraybackslash}p{0.19\textwidth}
  >{\raggedright\arraybackslash}p{0.19\textwidth}X}
\toprule
Target & H0 \(n=30\) & H1 \(n=25\) & Evaluator reading \\
\midrule
\(a\) & no solution & solved, two Horn rules & H0 desired; H1 spurious. \\
\(b\) & no solution & solved, symmetric rules & H0 desired; H1 spurious. \\
\(c\) & solved AND & solved AND & Control unchanged in kind. \\
\bottomrule
\end{tabularx}
\end{center}

\begin{lstlisting}
a(A) :- b_val_1(A), c_val_1(A).
a(A) :- b_val_0(A), c_val_0(A).
b(A) :- a_val_1(A), c_val_1(A).
b(A) :- a_val_0(A), c_val_0(A).
\end{lstlisting}

The both-zero rules are false population implications. They survive only
because no opposite-label witness occurs in those sample cells. H0 and H1
traces agree until the candidate is checked: H1 accepts it; H0 encounters
rows 26/30, attempts repair, and dead-ends.

\interpretation
This is support sensitivity, not a generic sample-size benefit. Five rows
matter because two instantiate the rare conflict. Nested prefixes control the
stream but are not independent replicates, so H1 gives no across-seed
frequency law.

\section{Irrelevant variables and BK order: H2 and H3}

\subsection{H2: trailing distractor under Greedy}

H2 adds isolated binary root \(D\) after the parents in target-\(c\) BK order.
Among \(n=30\) positive \(c\) rows, both \(D\) values occur nine times. AAMAS
emits:

\begin{lstlisting}
c(A) :- a_val_1(A), b_val_1(A), d_val_1(A).
c(A) :- a_val_1(A), b_val_1(A), d_val_0(A).
\end{lstlisting}

The pair is sample-adequate and jointly equivalent to AND on observed \(D\),
but is non-minimal and contains a causally irrelevant variable. Greedy
\code{folding_space(bk)} exhausts matching BK literals in each co-occurrence
body, so it retains \code{d_val_1} or \code{d_val_0}; no later cross-rule
compression removes \(D\).

The secondary \(n=2\) prefix yields only
\(\code{c :- a_val_1,b_val_1,d_val_1}\). This extreme
1-positive/1-negative case is not H2's centre.

\subsection{H3: leading distractor under nd}

\[
B\longrightarrow C\longleftarrow D,\qquad C=B\land D,
\]
with isolated \(A\) first in target-\(c\) order \(a,b,d\). The evaluator rule
is:

\begin{lstlisting}
c(A) :- b_val_1(A), d_val_1(A).
\end{lstlisting}

ECAI's first fold is:

\begin{lstlisting}
folding result: c(A) <- [a_val_1(A)]
\end{lstlisting}

and final target rules remain gated by the isolated variable:

\begin{lstlisting}
c(A) :- alpha_1(A), a_val_1(A).
c(A) :- alpha_2(A), a_val_0(A).
\end{lstlisting}

Contraries later mention \(b,d\), but the theory remains assumption-rich,
\(a\)-conditioned, and different from the evaluator BD conjunction. All 15
positives have \(b=d=1\), with both isolated \(a\) values present. The
first-fold choice reflects representation order, not causal relevance.

On the same table AAMAS emits:

\begin{lstlisting}
c(A) :- a_val_1(A), b_val_1(A), d_val_1(A).
c(A) :- a_val_0(A), b_val_1(A), d_val_1(A).
\end{lstlisting}

The mechanism appears in each body but is unnecessarily split by \(A\).

\subsection{H2--H3 synthesis}

\begin{itemize}
  \item Greedy maximal co-occurrence includes matching BK literals whether the
  isolated variable is trailing or leading.
  \item Nd is locally order-sensitive: a BK-leading irrelevant variable can
  become the first fold and organise later assumptions.
\end{itemize}

H2/H3 form a controlled fixture family, not a pure source-code ablation. Their
samples and identities differ. The trace explanations are supported, but do
not imply that every irrelevant variable is always retained.

\section{Learning semantics and search strategy: H4, H4b, and H5}

\subsection{H4: cautious semantics blocks the root residual gadget}

H4 fixes H0 fixture, sample, nd folding, selection, space, and \code{relto},
changing brave to repository-baseline cautious. Outcomes:

\begin{center}
\begin{tabular}{@{}lcc@{}}
\toprule
Target & Brave nd/\code{relto} & Cautious nd/\code{relto} \\
\midrule
\(a\) & solved & no solution \\
\(b\) & solved & no solution \\
\(c\) & solved & solved, byte-identical delta \\
\bottomrule
\end{tabular}
\end{center}

Root \(a\) traces coincide through early \(\alpha_1\)--\(\alpha_3\), then
diverge when closing \code{c_alpha_3} with \(\alpha_2,\code{b_val_0}\):

\begin{lstlisting}
% Brave:
found: alpha_2/1 ... OK
% Cautious:
found: alpha_2/1 ... KO
cannot introduce an assumption for [b_val_0(A)]
\end{lstlisting}

Brave completes the per-row gadget; cautious rejects it, explores further,
exhausts budget, and emits no delta. Target \(b\) mirrors this. Control \(c\)
still solves with the byte-identical ECAI delta: the semantic change rejects a
particular reuse, not every assumption-bearing solution.

\subsection{H4b: precise change on H3 target
\texorpdfstring{\(c\)}{c}}

Both brave and cautious nd/\code{relto} solve H3 \(c\). Their deltas differ in
exactly one rule:

\begin{lstlisting}
% Brave:
c_alpha_3(A) :- alpha_1(A), a_val_1(A).
% Cautious:
c_alpha_3(A) :- d_val_1(A).
\end{lstlisting}

Cautious rejects \(\alpha_1\) reuse on \code{a_val_1}, then accepts
\code{d_val_1}. Both arms retain the isolated-\(a\) target gates, the nested
\(\alpha_1\) contrary side, the flat \(\alpha_2\) side, and repair order through
\(\alpha_3\). Those shared structures arise from BK order, contrary-set
composition, and \code{relto} repair order, not cautious semantics.

\subsection{H5: the semantic flip is inert inside Greedy}

H5 duplicates five completed AAMAS collections with experimental
\code{greedy_cautious}. Across 18 paired cells:

\begin{itemize}
  \item all outcomes match;
  \item every solved delta is byte-identical;
  \item traces are isomorphic apart from entailment logging;
  \item there are no Greedy-cautious timeouts; and
  \item mean runtime is about \(1.84\times\), still below one second.
\end{itemize}

This includes H0 root no-solutions, H1 spurious roots, H2/H3 distractor
theories, and H0/H1 AND rules. Cautious semantics does not repair Greedy's
finite-sample or non-minimal outputs here.

\subsection{No isolated semantics ranking}

\begin{center}
\small
\begin{tabularx}{\textwidth}{
  >{\raggedright\arraybackslash}p{0.17\textwidth}
  >{\raggedright\arraybackslash}p{0.33\textwidth}X}
\toprule
Cell & Greedy-cautious & Baseline cautious nd \\
\midrule
H0 roots & no-solution in about \(0.4\)s & same outcome after about
\(5{,}000\) lines and \(66\)--\(70\)s \\
H0 \(c\) & direct Horn AND & assumption-rich \(a\)-gate, \(b=0\) contrary \\
H3 roots & short no-solution & timeout at 300s \\
H3 \(c\) & Horn rules split on isolated \(a\) & \(a\)-gated assumptions with
explicit \(d=1\) close \\
\bottomrule
\end{tabularx}
\end{center}

This is bundle-level: Greedy/mgr/BK and nd/any/all differ together. Cautious
semantics alone explains neither learned theory nor computational cost.

\section{Search-budget and sample-size cost: H6}

\subsection{H6a: folding budget}

At H0 \(n=30\), cautious-nd ceiling \(M\in\{1,2,5,10\}\) gives for target \(a\):

\begin{center}
\begin{tabular}{@{}rrrrr@{}}
\toprule
\(M\) & Outcome & Runtime (s) & Lines & New assumptions \\
\midrule
1 & no solution & 6.638 & 549 & 9 \\
2 & no solution & 13.117 & 1047 & 18 \\
5 & no solution & 32.775 & 2541 & 45 \\
10 & no solution & 69.952 & 5031 & 90 \\
\bottomrule
\end{tabular}
\end{center}

Target \(b\) mirrors this. Target \(c\) solves at every \(M\) with identical
delta, 134 lines, and about \(1.3\)s.

\code{folding_steps(M)} is a token ceiling, not exactly \(M\) successful folds.
The decisive root KO occurs in budget one; failure increments the allowance
and replays essentially the same motif:
\[
\operatorname{lines}(M)\approx C_{\mathrm{fixed}}+
M C_{\mathrm{failed\ band}}.
\]
There are roughly nine new assumptions and 500 lines per band. Through
\(M=10\), extra budget buys repeated failure, not a different explanation.

\subsection{H6b: nested sample size at fixed budget}

\begin{center}
\begin{tabular}{@{}rrrrr@{}}
\toprule
\(n\) & \(a\) lines & \(a\) runtime & \(c\) lines & \(c\) runtime \\
\midrule
30 & 1047 & 13.117 & 134 & 1.279 \\
60 & 1641 & 22.169 & 200 & 2.126 \\
90 & 2241 & 32.291 & 266 & 3.033 \\
\bottomrule
\end{tabular}
\end{center}

At fixed \(M=2\), root topology stays fixed: 18 new assumptions, 32 KOs, one
token increase. More rows enlarge grounded identifiers, rote rules,
subsumption checks, and contraries. All atoms already occur at \(n=30\);
multiplicities change from \(2/4/7/17\) to \(4/10/12/34\) to \(5/17/18/50\)
for \((0,0,0)/(0,1,0)/(1,0,0)/(1,1,1)\).

\interpretation
H4's long root traces are primarily repeated token bands, with secondary
per-band growth in \(n\). This is not a universal complexity law. No
\(n\times M\) grid or interaction result exists.

\section{Assumption introduction under nd: H7a and H7b}

\subsection{H7a design and roles}

H7a compares locked \code{ecai2024}/\code{relto} with experimental
\code{nd_brave_sechk}/\code{sechk}. Brave, nd, steps 10, selection, space,
post-fold checking, fixture, sample, BK, and examples are fixed in each pair.
H0 and H3 \(n=30\) give seven pairs.

\begin{center}
\begin{tabularx}{\textwidth}{
  >{\raggedright\arraybackslash}p{0.16\textwidth}
  >{\raggedright\arraybackslash}p{0.26\textwidth}X}
\toprule
Cell & Role & Evaluative use \\
\midrule
H0/H3 \(c\) & deterministic child & desirable mechanism target \\
H0 \(a\) & parent of \(c\), root as target & root diagnostic \\
H3 \(b\) & parent of \(c\), root as target & role-peer of H0 \(a\) \\
H3 \(a\) & isolated root & BK-leading distractor diagnostic \\
H0 \(b\), H3 \(d\) & other parents, roots as targets & supporting mirrors \\
\bottomrule
\end{tabularx}
\end{center}

\subsection{H7a seven-pair result}

Both arms solve every cell, but every paired delta differs:

\begin{center}
\small
\begin{tabularx}{\textwidth}{
  >{\raggedright\arraybackslash}p{0.10\textwidth}
  >{\raggedright\arraybackslash}p{0.14\textwidth}
  >{\raggedright\arraybackslash}p{0.20\textwidth}
  >{\raggedright\arraybackslash}p{0.20\textwidth}X}
\toprule
Cell & BK order & Outcome & \code{sechk} variable & Extra \code{relto} vars \\
\midrule
H0 \(a\) & \(b,c\) & both solved & \(b\) & \(c\) \\
H0 \(b\) & \(a,c\) & both solved & \(a\) & \(c\) \\
H0 \(c\) & \(a,b\) & both solved & \(a\) & \(b\) \\
H3 \(a\) & \(b,c,d\) & both solved & \(b\) & \(c,d\) \\
H3 \(b\) & \(a,c,d\) & both solved & \(a\) & \(c,d\) \\
H3 \(c\) & \(a,b,d\) & both solved & \(a\) & \(b,d\) \\
H3 \(d\) & \(a,b,c\) & both solved & \(a\) & \(b,c\) \\
\bottomrule
\end{tabularx}
\end{center}

For H0 \(c\):

\begin{lstlisting}
% relto
c(A) :- alpha_1(A), a_val_1(A).
c_alpha_1(A) :- b_val_0(A).
% sechk
c(A) :- alpha_1(A), a_val_1(A).
c_alpha_1(A) :- alpha_2(A), a_val_1(A).
c_alpha_2(A) :- alpha_1(A), a_val_1(A).
\end{lstlisting}

The outer \(a=1\) gate is shared. \code{relto} brings mechanism parent \(b\)
into the contrary; \code{sechk} closes on \(a\) alone. On H3 \(c\),
\code{relto} also brings \(b,d\), while \code{sechk} uses only isolated \(a\).

\subsection{H7a trace-supported explanation}

\begin{itemize}
  \item \code{relto} searches for an assumption relative to the folded body.
  A KO can backtrack to another fold literal, admitting later BK variables.
  \item \code{sechk} first mints a provisional assumption. Stable-extension
  checking may replace or retain it. Here it often closes before the
  alternative folds reached by \code{relto}.
\end{itemize}

\begin{lstlisting}
relto: gen2: ... looking for assumption relative to
sechk: gen2: generating NEW assumption: alpha_1(A)
\end{lstlisting}

Every H7a \code{sechk} delta cites the first predictor block and builds an
assumption--contrary chain on it: the bounded first-BK latch plus alpha-chain
pattern. It is not always \(a\); target \(a\)'s first predictor is \(b\).

Byte-identical H0/H3 \code{sechk} root deltas reflect a learner template, not
shared causal role. H0 \(a\) is a parent; H3 \(a\) is isolated; H3 \(b\) is
H0 \(a\)'s role-peer.

\boundary
H7a shows a brave-nd control-flow and theory difference. It does not rank
\code{relto}/\code{sechk}, prove a universal first-BK law, or show mechanism
recovery. Its first-BK latch is a bounded brave-nd pattern, not a prediction
about cautious learning.

\subsection{H7b design: the cautious counterpart}

H7b holds the H0/H3 fixtures, frozen \(n=30\) tables, exact-value encoding,
nd folding, steps 10, selection, space, and post-fold checking fixed. It
compares locked \code{baseline_cautious}/\code{relto} from H4/H4b with
experimental \code{nd_cautious_sechk}/\code{sechk}. The only
learning-relevant option changed within each pair is assumption introduction.
The comparator is not \code{ecai2024}: both H7b arms use cautious semantics.

Seven targets were run, but the scientific comparisons have unequal strength:

\begin{center}
\small
\begin{tabularx}{\textwidth}{
  >{\raggedright\arraybackslash}p{0.10\textwidth}
  >{\raggedright\arraybackslash}p{0.19\textwidth}
  >{\raggedright\arraybackslash}p{0.17\textwidth}X}
\toprule
Question & Cell & Outcome pair & Evidential use \\
\midrule
Q1 & H0 \(c\) & solved / solved & Comparable deltas and traces. \\
Q2 & H3 \(c\) & solved / solved & Comparable deltas and traces. \\
Q3 & H0 \(b\) & no solution / no solution & Failed-search cost only. \\
Q4 & H3 \(b\) & timeout / timeout & Non-comparable: both stdout files empty. \\
Q5 & H3 \(a\) & timeout / timeout & Non-comparable: both stdout files empty. \\
\bottomrule
\end{tabularx}
\end{center}

Only \(c\) is a desirable deterministic mechanism target. H0 \(b\) is a
parent/root diagnostic; H3 \(b\) is its role-peer, whereas H3 \(a\) is an
isolated distractor. Matching timeout labels on Q4--Q5 do not establish
matching search because no trace was serialized.

\subsection{H7b Q1--Q2: cautious child theories diverge}

Both cautious policies solve \(c\) on both tables, but the paired deltas are
different. Learner-visible data, BK, and examples are byte-identical within
each pair.

For Q1, the H0 \(c\) theories are:

\begin{lstlisting}
% baseline_cautious / relto
c(A) :- alpha_1(A), a_val_1(A).
c_alpha_1(A) :- b_val_0(A).

% nd_cautious_sechk / sechk
c(A) :- alpha_1(A), a_val_1(A).
c_alpha_1(A) :- alpha_2(A), a_val_1(A).
c_alpha_2(A) :- b_val_1(A).
\end{lstlisting}

The outer \(a=1\) gate is shared. Relative introduction closes the contrary
directly with \code{b_val_0}; \code{sechk} first inserts another assumption
layer on \code{a_val_1}, then reaches \code{b_val_1}. It therefore does not
ignore parent \(b\), but represents the repair at greater depth.

For Q2, both H3 \(c\) theories retain the isolated-\(a\) gates and cite the
mechanism parents \(b\) and \(d\) somewhere in the contrary layer:

\begin{center}
\begin{tabularx}{\textwidth}{
  >{\raggedright\arraybackslash}p{0.24\textwidth}
  >{\raggedright\arraybackslash}p{0.18\textwidth}
  >{\raggedright\arraybackslash}p{0.17\textwidth}X}
\toprule
Arm & Assumptions & Trace lines & Character \\
\midrule
H4b \code{relto} & 3 & 257 & Lean relative-KO path to \(b,d\). \\
H7b \code{sechk} & 5 & 899 & Deeper nests and repeated NEW/\code{gen5}/KO cycles. \\
\bottomrule
\end{tabularx}
\end{center}

The thicker \code{sechk} theory is still distractor-retaining because its
ordinary target rules remain gated by isolated \(a\). Retaining \(b,d\) in
contraries is therefore not mechanism recovery.

\subsection{H7b trace explanation}

H7a and H7b enter the same implementation fork after failed post-fold
entailment:

\begin{lstlisting}
relto: gen2: ... looking for assumption relative to
sechk: gen2: generating NEW assumption: alpha_1(A)
\end{lstlisting}

Under \code{relto}, a relative assumption is tried first. Under cautious
\code{gen3}, reuse often receives KO, permitting backtracking to another fold
literal. Under \code{sechk}, a provisional assumption is always minted and
tested through \code{gen5}/\code{gen6}; cautious rejection can then force
further search and eventually expose later \(b\) or \(d\) folds. Thus the
control-flow fork is shared with H7a, but cautious entailment changes which
repairs survive and how far search proceeds.

This explains the bounded contrast: brave \code{sechk} closes early around
the first predictor on the two child cells, whereas cautious \code{sechk}
retains mechanism-parent covariates through later, deeper repairs. It is not
evidence that cautious \code{sechk} is generally better.

\subsection{H7b Q3--Q5: cost and evidence limits}

For H0 target \(b\), neither cautious arm emits a delta. Both traverse the
same folding-token ladder and finish with no solution, but \code{sechk}
produces approximately 13,818 trace lines versus 5,118 for \code{relto}, with
many more NEW/\code{gen5} attempts. This is a failed-search cost difference,
not a learned-theory comparison and not root mechanism recovery.

For H3 targets \(b\) and \(a\), both arms time out at 300 seconds and both
stdout files contain zero bytes. Q4--Q5 therefore say only that the configured
processes timed out; they provide no basis for saying that the search paths or
hypotheses matched. The H3 target \(d\) timeout is supporting collection
context, not one of the predeclared Q1--Q5 comparisons.

\subsection{H7a--H7b synthesis}

\begin{center}
\small
\begin{tabularx}{\textwidth}{
  >{\raggedright\arraybackslash}p{0.24\textwidth}
  >{\raggedright\arraybackslash}p{0.31\textwidth}X}
\toprule
Property & H7a: brave nd & H7b: cautious nd \\
\midrule
Within-mode intervention & \code{relto} versus \code{sechk}; seven comparable
solved deltas & \code{relto} versus \code{sechk}; two comparable solved child
deltas, one failed-cost pair, two empty-log timeout pairs \\
Child covariates under \code{sechk} & First-BK latch; later mechanism parents
absent & Mechanism parents retained through deeper contrary nests \\
Search character & Early assumption-chain closure & More cautious KO,
backtracking, nesting, and cost \\
Scientific reading & Assumption policy changes theory structure & The effect
of that policy is conditioned by entailment semantics \\
\bottomrule
\end{tabularx}
\end{center}

Together the probes show that assumption introduction is a genuine procedural
intervention, but its observable effect is not invariant to brave versus
cautious evaluation. This is a bounded two-fixture interaction pattern, not a
universal law or method ranking.

\section{Cross-cutting synthesis}

\subsection{What is recoverable}

The deterministic \(c\) mechanism is representable in exact-value BK. H0 shows
Greedy can emit it when positives share one parent configuration and no
irrelevant predictors exist. This is a positive rule-level capability.

It is not robust across interventions:
\begin{itemize}
  \item nd/\code{relto} represents the same relation with assumption/contrary;
  \item Greedy retains isolated variables in H2/H3;
  \item nd makes a leading isolated variable the outer gate in H3;
  \item brave-nd \code{sechk} drops later mechanism parents in H7a;
  \item cautious-nd \code{sechk} retains those parents in H7b, but through
  deeper distractor-gated theories rather than the evaluator conjunction; and
  \item no target-wise outputs are automatically aggregated into a graph.
\end{itemize}
The evidence supports configuration-specific mechanism-rule recovery, not
target-agnostic graph recovery.

\subsection{\code{solved} versus correctness}

\begin{enumerate}
  \item H0 AAMAS \(c\): evaluator mechanism string.
  \item H0 ECAI \(c\): defeasible, non-canonical sample representation.
  \item H2/H3 \(c\): irrelevant-variable dependence.
  \item H7b cautious \(c\): solved with deeper, non-canonical
  assumption/contrary theories under \code{sechk}.
  \item Brave roots: assumption solutions despite no deterministic mechanism.
\end{enumerate}
Outcome, semantic adequacy, syntactic alignment, minimality, and causal
recovery must remain separate.

\subsection{Failure attribution}

\begin{center}
\begin{tabularx}{\textwidth}{
  >{\raggedright\arraybackslash}p{0.23\textwidth}
  >{\raggedright\arraybackslash}p{0.25\textwidth}X}
\toprule
Source & Evidence & Explanation \\
\midrule
Finite support & H1 & Rare counterexample absent, so false root rule survives. \\
Representation/order & H2, H3 & Isolated variables enter maximal bodies or
the first nd fold. \\
Target geometry & H0 & Child has one positive parent cell; roots have
opposite labels within predictor cells. \\
Search strategy & H0, H3, H5 & Greedy and nd build different theories and
incur different costs. \\
Learning semantics & H4, H4b, H5 & Brave reuse may be accepted where cautious
rejects it; effect depends on bundle. \\
Search budget & H6a & Larger token ceiling repeats a failed motif. \\
Grounding size & H6b & More rows enlarge fixed-topology bookkeeping. \\
Assumption policy and semantics & H7a, H7b & \code{relto}/\code{sechk} take
different repair paths; which covariates survive and at what cost depends on
brave versus cautious checking. \\
Evaluation objective & All & Example satisfaction does not require minimality,
functional determination, or graph recovery. \\
\bottomrule
\end{tabularx}
\end{center}

\subsection{Strategy failure versus observational non-identifiability}

The baseline population is ordinarily faithful to a singleton-MEC collider.
An ideal infinite-data CI oracle would identify its ordinary CPDAG. The
target-wise learner receives no CI tests, acyclicity theory, or graph objective,
so non-canonical rules cannot be attributed automatically to Markov
equivalence.

Root ambiguity is instead an information limitation of the encoded
predictor-to-target task: identical BK configurations have both labels. Brave
semantics can construct a witness, but that does not make the root a function;
H4 shows cautious semantics can refuse the construction.

\begin{itemize}
  \item H1: finite-sample support.
  \item H2/H3: representation and strategy.
  \item H4/H4b/H5: semantics conditioned by search.
  \item H6: procedural cost.
  \item H7a/H7b: assumption-introduction control flow conditioned by
  entailment semantics, with timeout-limited comparisons on H3 roots.
\end{itemize}
None is by itself a general observational non-identifiability result.

\subsection{Contribution and missing machinery}

ABA Learning contributes interpretable rules, assumptions, contraries, and
traces showing how conflicts are handled. The current bridge lacks:

\begin{itemize}
  \item aggregation from target-wise frameworks to DAG/CPDAG;
  \item causal constraints distinguishing parents, proxies, and distractors;
  \item invariance or minimality across rule presentations;
  \item a general semantic-equivalence test against the mechanism; and
  \item protection against finite-sample support gaps.
\end{itemize}

Later Causal ABA machinery might constrain graph-compatible hypotheses, expose
independence information, guide folds, or prefer simpler compatible
frameworks. These are motivations, not H0--H7b results.

\section{Candidate claims for supervisor discussion}

No claim below is locked. Each is scoped to completed evidence and should be
reviewed with Fabrizio before any claim is approved.

\subsection*{C1: procedural solution is insufficient}
\textbf{Statement.} On these deterministic colliders, \code{solved} does not
establish mechanism recovery: solved theories may be canonical,
assumption-mediated, distractor-retaining, or non-functional under brave
grounding.

\evidence H0 contrasts child conjunction and root choice; H2/H3 solve \(c\)
with isolated variables; H7a solves all cells with different deltas; H7b
solves both child cells with further non-canonical deltas.

\boundary This concerns insufficiency of the label, not uselessness of solved
theories or impossibility of semantic equivalence.

\subsection*{C2: finite support can create spurious root rules}
\textbf{Statement.} Under H0 Greedy AAMAS, omitting rare \((0,0,0)\) from the
\(n=25\) prefix permits false both-zero root rules; its \(n=30\) presence
blocks them.

\evidence H1 fixes fixture, seed sequence, encoding, and configuration; traces
diverge at rows 26/30.

\boundary One seed/prefix contrast; not a general sample-complexity result.

\subsection*{C3: irrelevant-variable retention is strategy-specific}
\textbf{Statement.} On tested four-variable AND fixtures, Greedy and nd retain
an isolated predictor by different procedures: maximal co-occurrence versus
BK-leading first-fold repair.

\evidence H2 supplies trailing-\(D\) Greedy; H3 supplies leading-\(A\) nd and
an AAMAS contrast.

\boundary Bounded fixtures/encodings, not a complete permutation study.

\subsection*{C4: semantic effects depend on search}
\textbf{Statement.} Brave-to-cautious changes learned behaviour under
nd/\code{relto} but not outcomes or solved deltas under Greedy/\code{relto} on
the tested collections.

\evidence H4 changes root outcomes; H4b changes one \(c\) contrary closure; H5
has 18/18 outcome matches and byte-identical solved deltas; H7a/H7b show that
the \code{sechk} effect itself changes across brave and cautious contexts.

\boundary No global brave/cautious ranking; direct nd/Greedy comparison is
multi-option.

\subsection*{C5: cautious-nd root cost is replayed failure}
\textbf{Statement.} On H0 roots, cost grows with folding-token ceiling because
budgets replay a similar failed motif; at fixed budget, nested \(n\) enlarges
per-band work.

\evidence H6a has about nine new assumptions per budget and near-linear line
growth; H6b fixes NEW/KO topology while lines grow for \(n=30,60,90\).

\boundary Limited slices, not universal complexity or an \(n\times M\)
interaction.

\subsection*{C6: assumption policy changes covariate structure}
\textbf{Statement.} Under nd on the tested H0/H3 tables, switching
\code{relto} to \code{sechk} changes learned child theories, but the form of
that change depends on entailment semantics. Brave \code{sechk} follows a
first-BK chain and omits later mechanism parents; cautious \code{sechk} can
reach those parents through deeper assumption nests and increased search.

\evidence Seven exact input-matched brave H7a pairs have different deltas;
both exact input-matched cautious H7b child pairs also have different deltas.
Their traces share the same initial \code{gen.pl} fork, while cautious checks
change which repairs survive. H7b Q3 additionally shows greater failed-search
work under \code{sechk}.

\boundary Two fixtures, one seed, and only two completed cautious child-theory
comparisons. H7b Q4--Q5 are non-comparable empty-log timeouts. No method
ranking, universal interaction law, or mechanism-recovery conclusion.

\section{Limitations and questions for Fabrizio}

\subsection{Limitations}
\begin{itemize}
  \item Narrow family: binary stochastic roots and deterministic AND child.
  \item Mostly one seed and principal \(n=30\); prefixes are not replications.
  \item H2/H3 cover two distractor positions, not all orders/mechanisms.
  \item Binary exact-value encoding only; no missing/non-discrete data here.
  \item Ordinary faithfulness does not settle all deterministic identifiability.
  \item No general semantic equivalence checker for learned theories.
  \item No target-union to graph, MEC, or CPDAG.
  \item H4b non-\(c\) cautious cells timed out at 300s and were out of scope.
  \item H7b Q4--Q5 timed out with empty traces, preventing comparison of the
  H3 root search paths.
  \item The brave/cautious assumption-policy comparison is not complete for
  every target because cautious H3 roots did not produce inspectable traces.
\end{itemize}

\subsection{Questions}
\begin{enumerate}
  \item Should semantically equivalent defeasible theories count as mechanism
  recovery without syntactic identity?
  \item What minimality/invariance condition should reject H2/H3 distractors?
  \item Is root \code{completed_no_solution} a guardrail, non-result, or both?
  \item What graph object should a target-wise collection produce?
  \item Which causal constraints should later become learner-visible?
  \item Does the semantics-conditioned \code{relto}/\code{sechk} pattern
  persist across another mechanism, sample, or BK order?
  \item Which candidate claims need another fixture or seed before approval?
\end{enumerate}

\section{Conclusion}

H0--H7b provide a controlled, trace-backed account of target-wise ABA Learning
on deterministic causal fixtures. Greedy can recover the direct AND rule when
the representation is minimal and relevant sample cells are present. That
behaviour is not stable across support, irrelevant predictors, search order,
semantics, or assumption policy.

The probes locate rather than merely enumerate failures: H1 identifies
support; H2/H3 representation/order; H4/H5 semantics versus search; H6 cost;
H7a/H7b assumption-introduction control flow and its dependence on entailment
semantics. This attribution is the principal contribution of the investigation.

The evidence does not show causal graph discovery. It does show where unguided
symbolic learning produces a mechanism-aligned rule, where it instead builds
predictive or defeasible alternatives, and what later causal guidance may
need to constrain.

\appendix

\section{Compact evidence index}

The records below share this directory:

\begin{lstlisting}
docs/experiments/qualitative/M13-C3-binary-collider-and/
\end{lstlisting}

Learner artefacts are relative to
\path{causal/outputs/aba_learning/targetwise/}.

\begin{itemize}
  \item \textbf{H0:} \path{learning_analysis.md}; H0 fixture collections
  \path{aamas2025/n30_seed42} and \path{ecai2024/n30_seed42}.

  \item \textbf{H1:} \path{h1_support_ablation.md}; H0 fixture collection
  \path{aamas2025/n25_seed42}.

  \item \textbf{H2:} \path{h2_irrelevant_covariate.md}; collections below
  \path{m13_bucket3_binary_collider_and_iso_d/aamas2025/}.

  \item \textbf{H3:} \path{h3_bk_leading_distractor.md}; collections below
  \path{m13_bucket3_binary_bd_and_lead_a/} for \path{ecai2024} and
  \path{aamas2025} at \path{n30_seed42}.

  \item \textbf{H4:} \path{h4_cautious_vs_brave.md}; H0 fixture collection
  \path{baseline_cautious/n30_seed42}.

  \item \textbf{H4b:} \path{h4b_cautious_split_under_a.md}; H3 fixture
  collection \path{baseline_cautious/n30_seed42}.

  \item \textbf{H5:} \path{h5_greedy_cautious.md}; all five completed AAMAS
  fixture/sample counterparts under \path{greedy_cautious}.

  \item \textbf{H6:} \path{h6_folding_and_n_ablation.md}; H0 fixture
  collections under \path{baseline_cautious_steps1},
  \path{baseline_cautious_steps2}, \path{baseline_cautious_steps5}, and the
  reused \path{baseline_cautious} steps-10 collection.

  \item \textbf{H7a:} \path{h7a_relto_vs_sechk.md}; H0 and H3 fixture
  collections under \path{nd_brave_sechk/n30_seed42}.

  \item \textbf{H7b:} \path{h7b_cautious_relto_vs_sechk.md}; H0 and H3
  fixture collections under \path{nd_cautious_sechk/n30_seed42}, compared
  with the locked \path{baseline_cautious/n30_seed42} collections.
\end{itemize}

\section{Selected provenance}
\begin{itemize}
  \item H0 spec: \path{causal/fixtures/specs/m13_bucket3_binary_collider_and.yaml}.
  \item H0 reference:
  \path{causal/outputs/causal_fixtures/m13_bucket3_binary_collider_and/mechanism_reference.json}.
  \item H3 reference:
  \path{causal/outputs/causal_fixtures/m13_bucket3_binary_bd_and_lead_a/mechanism_reference.json}.
\end{itemize}

Frozen sample SHA-256 digests:

\begin{lstlisting}[basicstyle=\ttfamily\footnotesize]
H0 n30_seed42
1edae465d0f6b6b662b3c4ed4863af21c1d9b0c1cf428774731dd3f43282a8ca
H1 n25_seed42
c7ccf94c10a7862f34b7eeedb61529cf60ac0b13b932ac608beaea1a6d17fdd4
H3 n30_seed42
48bd0da0df696e29fb91e16e79bf30592b55c6dc16f34ee72e74aa36cc52db05
\end{lstlisting}

\section{Reproducibility pattern}

\begin{lstlisting}
python -m causal.targetwise.cli validate --config <targetwise-config.yaml>
python -m causal.targetwise.cli run --config <targetwise-config.yaml>
\end{lstlisting}

Principal per-cell evidence:
\begin{itemize}
  \item \path{input/data.csv}: frozen learner-visible table;
  \item \path{input/bk.aba}: exact-value BK;
  \item \path{input/examples.json}: target-specific examples;
  \item \path{output/prolog.stdout}: transformation/search trace;
  \item \path{output/delta.aba}: learned additions when solved;
  \item \path{output/bk.sol.aba}, \path{output/bk.sol.asp}, and
  \path{output/bk.sol_chk.asp}: serialized and checked artefacts;
  \item \path{metrics.json}: outcome/runtime/rule diagnostics; and
  \item \path{summary.md}: collection outcome index.
\end{itemize}

Fixture populations, certificates, and \path{mechanism_reference.json} remain
evaluator-only and are not copied into learner BK/examples.

\section{Claim-status reminder}

Candidate claims are proposals for discussion with Fabrizio. The canonical
Bucket~3 hub still records no approved claim. H7a and H7b are complete at the
probe level; their cross-semantic interpretation remains bounded by the H7b
timeouts. Later report wording must be checked against canonical experiment
records and the claims ledger.

\end{document}
